\clearpage\thispagestyle{empty}
\pdfbookmark[0]{Quellen zum Pallas-Bericht vom Dezember 1802}{pallas180212-sources}
\begin{center}\textbf{Quellen zum Pallas-Bericht vom Dezember 1802}\end{center}
\noindent Diese KI-Übertragung in heutiges Deutsch folgt dem Auszug aus Franz Xaver von Zachs Pallas-Bericht vom Dezember 1802 in \href{https://gdz.sub.uni-goettingen.de/id/PPN236007467}{Gauß’ \textit{Werken}, Band~VI} (Göttingen, 1874), S.~232--234. Er enthält Gauß’ Vergleich mit Orianis Beobachtungen, die fünfte Fassung der Elemente, den Vergleich mit Messiers letzter Beobachtung und die Angaben zur Lichtstärke.
\par\smallskip\noindent Der \href{https://api.digitale-sammlungen.de/iiif/image/v2/bsb10538598_00593/full/1450,/0/default.jpg}{Erstdruck} steht als Beitrag LXV in der \textit{Monatlichen Correspondenz}, Band~VI, S.~579--583. Der Auszug beginnt mit dem \href{https://api.digitale-sammlungen.de/iiif/image/v2/bsb10538598_00594/full/1450,/0/default.jpg}{Vergleich auf S.~580} und setzt sich mit den \href{https://api.digitale-sammlungen.de/iiif/image/v2/bsb10538598_00595/full/1450,/0/default.jpg}{Elementen V auf S.~581} und dem \href{https://api.digitale-sammlungen.de/iiif/image/v2/bsb10538598_00596/full/1450,/0/default.jpg}{Vergleich mit Messier auf S.~582} fort. Die Zeitschrift druckt auf \href{https://api.digitale-sammlungen.de/iiif/image/v2/bsb10538598_00597/full/1450,/0/default.jpg}{S.~583 die vollständige Ephemeride und vier Lichtstärken}. Die \textit{Werke} beschreiben die Ephemeride in eckigen Klammern und geben eine Handschriftenlesart zur Lichtstärke im April und vier weitere Lichtstärken an.
\par\smallskip\noindent Der Bericht verweist auf die \href{https://api.digitale-sammlungen.de/iiif/image/v2/bsb10538598_00404/full/1450,/0/default.jpg}{Elemente IV} und die \href{https://api.digitale-sammlungen.de/iiif/image/v2/bsb10538598_00405/full/1450,/0/default.jpg}{frühere Ephemeride} im Oktoberheft und auf \href{https://api.digitale-sammlungen.de/iiif/image/v2/bsb10538598_00516/full/1450,/0/default.jpg}{Messiers Beobachtung} im Novemberheft.
\par\smallskip\noindent Übertragung und Satz wurden mit KI erstellt. Die Digitalisate der Zeitschrift stammen von der Bayerischen Staatsbibliothek.
\clearpage
